\documentclass[10 pt,usenames,dvipsnames, oneside]{article}
\usepackage{../../../modelo-ensino-medio}



\begin{document}

\begin{center}
  \begin{minipage}[l]{3cm}
\includegraphics[width=2cm]{logo}    
\end{minipage}\hfill
\begin{minipage}[r]{.8\textwidth}
 {\Large \scshape Atividade: Relações entre os logaritmos}  
\end{minipage}
\end{center}
\vspace{.2cm}

\ifdefined\prof
%Habilidades da BNCC
\begin{objetivos}
\item \textbf{EM13MAT305} Resolver e elaborar problemas com funções logarítmicas nos quais seja necessário compreender e interpretar a variação das grandezas envolvidas, em contextos como os de abalos sísmicos, pH, radioatividade, Matemática Financeira, entre outros.
\end{objetivos}

%Caixa do Para o Professor
\begin{goals}
%Objetivos específicos
\begin{enumerate}
\item Observar padrões nos valores dos logaritmos quando os logaritmandos são multiplicados por um número natural.
\end{enumerate}

\end{goals}

\bigskip
\begin{center}
{\large \scshape Atividade}
\end{center}
\fi

Utilize a \hyperref[potencias2_log]{tabela \ref{potencias2_log}} para encontrar os valores (aproximados) dos seguintes logaritmos em seu caderno:
\begin{enumerate}
\item $\log_2 2$, $\log_2 4$, $\log_2 8$ e $\log_2 16$.
\item $\log_2 3$, $\log_2 6$, $\log_2 12$ e $\log_2 24$.
\item $\log_2 3$, $\log_2 9$ e $\log_2 27$.
\item $\log_2 3$, $\log_2 5$ e $\log_2 15$.
\end{enumerate}

\begin{table}[H]
\centering
\setlength\tabulinesep{1.5pt}
\begin{tabu} to \textwidth{|>{$}l<{$}|>{$}c<{$}|}
\hline
\rowfont{\color{white}}\rowcolor{\currentcolor!80}
\bm{2^{\log_2n}} & \bm{n} \tabularnewline
\hline
2^0 & 1 \tabularnewline
\hline
2^1 & 2 \tabularnewline
\hline
2^{1{,}58} & 3 \tabularnewline
\hline
2^{2} & 4 \tabularnewline
\hline
2^{2{,}32} & 5 \tabularnewline
\hline
2^{2{,}58} & 6 \tabularnewline
\hline
2^{2{,}81} & 7 \tabularnewline
\hline
2^{3} & 8 \tabularnewline
\hline
2^{3{,}16} & 9 \tabularnewline
\hline
2^{3{,}32} & 10 \tabularnewline
\hline
\end{tabu}\hspace{2em}
\begin{tabu} to \textwidth{|>{$}l<{$}|>{$}c<{$}|}
\hline
\rowfont{\color{white}}\rowcolor{\currentcolor!80}
\bm{2^{\log_2n}} & \bm{n} \tabularnewline
\hline
2^{3{,}46} & 11 \tabularnewline
\hline
2^{3{,}58} & 12 \tabularnewline
\hline
2^{3{,}7} & 13 \tabularnewline
\hline
2^{3{,}81} & 14 \tabularnewline
\hline
2^{3{,}9} & 15 \tabularnewline
\hline
2^{4} & 16 \tabularnewline
\hline
2^{4{,}08} & 17 \tabularnewline
\hline
2^{4{,}16} & 18 \tabularnewline
\hline
2^{4{,}24} & 19 \tabularnewline
\hline
2^{4{,}32} & 20 \tabularnewline
\hline
\end{tabu}\hspace{2em}
\begin{tabu} to \textwidth{|>{$}l<{$}|>{$}c<{$}|}
\hline
\rowfont{\color{white}}\rowcolor{\currentcolor!80}
\bm{2^{\log_2n}} & \bm{n} \tabularnewline
\hline
2^{4{,}39} & 21 \tabularnewline
\hline
2^{4{,}46} & 22 \tabularnewline
\hline
2^{4{,}52} & 23 \tabularnewline
\hline
2^{4{,}58}& 24 \tabularnewline
\hline
2^{4{,}64} & 25 \tabularnewline
\hline
2^{4{,}7} & 26 \tabularnewline
\hline
2^{4{,}75} & 27 \tabularnewline
\hline
2^{4{,}81} & 28 \tabularnewline
\hline
2^{4{,}85} & 29 \tabularnewline
\hline
2^{4{,}9} & 30 \tabularnewline
\hline
\end{tabu}

\caption{Expoentes de $2$ aproximando os naturais de $1$ à $30$}\label{potencias2_log}
\end{table}


\ifdefined\prof
\begin{solucao}

	\begin{enumerate}
	\item $\log_2 2 = 1$, $\log_2 4 = 2$, $\log_2 8 = 3$ e $\log_2 16 = 4$.
	\item $\log_2 3 \approx 1{,}58$, $\log_2 6 \approx 2{,}58$, $\log_2 12 \approx 3{,}58$ e $\log_2 24 \approx 4{,}58$.
	\item $\log_3 3 = 1$, $\log_3 9 = 2$ e $\log_3 27 = 3$.
	\item $\log_2 3 \approx 1{,}58$, $\log_2 5 \approx 2{,}32$ e $\log_2 15 \approx 3{,}9$.
	\end{enumerate}

\end{solucao}
\fi

\end{document}